% ============================================================
% TESIS DOCTORAL
% Eficiencia hospitalaria, estructura organizativa
% y teoría de agencia: Un análisis de frontera estocástica
% del sistema hospitalario español
%
% José María Elena Izquierdo
% Departamento de Economía Aplicada
% Universidad de Salamanca
% 2026
% ============================================================

\documentclass[12pt, a4paper, twoside, openright]{book}

% ── Codificación y lengua ───────────────────────────────────
\usepackage[utf8]{inputenc}
\usepackage[T1]{fontenc}
\usepackage[spanish,es-nodecimaldot,es-noquoting]{babel}

% ── Tipografía ──────────────────────────────────────────────
\usepackage{lmodern}
\usepackage{microtype}

% ── Matemáticas ─────────────────────────────────────────────
\usepackage{amsmath, amssymb, amsthm, amsfonts}
\usepackage{bm}

% ── Tablas y figuras ────────────────────────────────────────
\usepackage{booktabs}
\usepackage{longtable}
\usepackage{array}
\usepackage{multirow}
\usepackage{threeparttable}
\usepackage{tabularx}
\usepackage{rotating}
\usepackage{pdflscape}
\usepackage{graphicx}
\usepackage{caption}
\usepackage{subcaption}
\usepackage{float}

% ── Diseño de página ────────────────────────────────────────
\usepackage{geometry}
\geometry{
  top=2.5cm, bottom=2.5cm,
  left=3.5cm, right=2.5cm,
  bindingoffset=0.5cm,
  headheight=28pt
}
\usepackage{setspace}
\onehalfspacing

% ── Encabezados y pies ──────────────────────────────────────
\usepackage{fancyhdr}
\pagestyle{fancy}
\fancyhf{}
\fancyhead[LE]{\small\itshape\leftmark}
\fancyhead[RO]{\small\itshape\rightmark}
\fancyfoot[C]{\thepage}
\renewcommand{\headrulewidth}{0.4pt}

% ── Color (necesario para hyperref con colores mixtos) ──────
\usepackage[dvipsnames,svgnames,x11names]{xcolor}

% ── Referencias y enlaces ──────────────────────────────────
\usepackage[colorlinks=true,
  linkcolor={blue!60!black},
  citecolor={blue!60!black},
  urlcolor={blue!60!black}]{hyperref}
\usepackage{natbib}
\usepackage{eurosym}

% ── Figuras: ruta hacia directorio teórico ─────────────────
\graphicspath{{../Teórico/}{./}}

% ── Diagramas TikZ ──────────────────────────────────────────
\usepackage{tikz}
\usetikzlibrary{arrows.meta, positioning, shapes.geometric, babel}

% ── Otros ───────────────────────────────────────────────────
\usepackage{color,soul}
\usepackage{enumitem}
\usepackage{footnote}

% ── Comandos propios ────────────────────────────────────────
\newcommand{\E}{\mathbb{E}}
\newcommand{\CE}{\mathrm{CE}}
\newcommand{\N}{\mathcal{N}}
\newcommand{\e}{\mathrm{e}}
\newcommand{\bsym}[1]{\boldsymbol{#1}}
\newcommand{\sig}[1]{\ensuremath{^{#1}}}

% ── Entornos de teorema ─────────────────────────────────────
\theoremstyle{plain}
\newtheorem{proposition}{Proposición}[chapter]
\newtheorem{result}{Resultado}[chapter]
\newtheorem{corollary}{Corolario}[chapter]
\theoremstyle{definition}
\newtheorem{definition}{Definición}[chapter]
\theoremstyle{remark}
\newtheorem{remark}{Observación}[chapter]

% Estilo compartido para ambos cronogramas
\tikzset{
  cronograma/.style={
    >=Stealth, thick
  },
  cdot/.style={
    circle, fill=black, minimum size=5pt, inner sep=0pt
  },
  clbl/.style={
    font=\fontsize{8.5}{10}\selectfont, align=center, 
    text width=3.2cm
  }
}% ═══════════════════════════════════════════════════════════
